% !TeX root = ../../tesis.tex

% =============================================================================
%  B.7 — GEOLAYERS: SERIALIZACIÓN PARA TRANSPORT-GIS
% =============================================================================
\section{GeoLayers: Serialización para Transport-\textsc{gis}}
\label{anx:vft-geolayers}

El módulo \path{src/api/routes/geo_layers.py} es la capa de salida geoespacial
de VFTModel. No contiene lógica analítica: recibe los \texttt{DataFrame} de
los algoritmos topológicos y los serializa como \texttt{FeatureCollection}
en \textsc{crs} \textsc{epsg:4326} (WGS84), consumibles directamente por
Transport-\textsc{gis}-\textsc{zmvm} y \textsc{qgis}.

Cada respuesta \texttt{FeatureCollection} lleva un bloque de metadatos
estándar que incluye el nombre del indicador, el número de geometrías, el
sistema de referencia y la marca temporal de la consulta.

\begin{lstlisting}[style=estiloPython,
    caption={Envoltura estándar VFT para toda respuesta GeoJSON.
             \texttt{src/api/routes/geo\_layers.py}.},
    label=lst:vft-feature-collection]
def _build_feature_collection(indicador, layer, features, parametros):
    return {
        "type": "FeatureCollection",
        "metadata": {
            "indicador":  indicador,
            "layer":      layer,
            "n_features": len(features),
            "crs":        "EPSG:4326",
            "fetched_at": datetime.now(timezone.utc).strftime("%Y-%m-%dT%H:%M:%SZ"),
            "parametros": parametros
        },
        "features": features
    }
\end{lstlisting}

\subsection*{Clasificadores de umbrales cualitativos}

Para facilitar la visualización en Transport-\textsc{gis}, cada entidad
lleva un campo categórico derivado de su valor numérico:

\begin{lstlisting}[style=estiloPython,
    caption={Clasificadores cualitativos para cobertura, fuerza capilar y
             factor de desviación. \texttt{src/api/routes/geo\_layers.py}.},
    label=lst:vft-clasifiers]
def _clasify_coverage(pct: float) -> str:
    if pct > 60:  return "alta"
    if pct >= 30: return "media"
    return "baja"

def _clasify_node(fc_total: float) -> str:
    if fc_total > 20: return "hub_principal"
    if fc_total > 10: return "hub_secundario"
    if fc_total > 6:  return "nodo_relevante"
    return "nodo_basico"

def _clasify_diff(factor: float) -> str:
    if factor <= 1.3: return "eficiente"
    if factor <= 1.6: return "moderado"
    if factor <= 2.0: return "alto"
    return "critico"
\end{lstlisting}

\subsection*{Capas disponibles en \texttt{/geolayers/coverage}}

El endpoint de cobertura expone tres subtipos mediante el parámetro
\texttt{layer}: \texttt{cobertura\_por\_alcaldia} (polígono de cada demarcación
coloreado por porcentaje de cobertura), \texttt{estaciones} (puntos de toda la
red sin duplicados) y \texttt{cobertura\_800m} (mancha vectorial de la unión
de buffers de 800\,m alrededor de cada parada). Las geometrías internas se
procesan en \textsc{epsg:32614} para operaciones métricas precisas y se
reprojectan a \textsc{epsg:4326} antes de la serialización.
